\section{Por qué Coding se parece a una GPU: un acelerador, no una aplicación}

Comparar el modelo de Coding con una GPU es la analogía con mayor poder explicativo de este episodio. La GPU no es en sí una aplicación de consumo, pero hace posibles el entrenamiento y la inferencia; del mismo modo, el modelo de Coding no es solo un complemento del IDE, sino una infraestructura que acelera la investigación, la ingeniería, el procesamiento de datos, los experimentos y la iteración de productos.

\subsection{Tres vías de aceleración}

\noindent\begin{tabularx}{\textwidth}{p{0.22\textwidth}p{0.34\textwidth}X}
\toprule
Vía de aceleración & Dónde ocurre & Efecto \\
\midrule
Amplificación de la productividad individual & Investigadores, programadores y arquitectos de primer nivel & El 1\% del talento se amplifica de 10 a 50 veces y se acorta el ciclo que va de la idea al experimento. \\
Capacidad de producción de I+D de la organización & feature, pipeline de datos, experimentos multimodales & Iteraciones que antes tomaban semanas se reducen a días y aumenta la frecuencia de lanzamiento de productos. \\
IA que acelera la IA & El modelo ayuda a escribir código de entrenamiento, scripts de evaluación y a depurar experimentos & La propia investigación en IA se acelera con herramientas de IA, lo que forma un ciclo de retroalimentación recursivo. \\
\bottomrule
\end{tabularx}

\begin{knowledgebox}{Por qué el código es la palanca del mundo digital}
El lenguaje natural expresa problemas; el código expresa soluciones. Siempre que una tarea ocurra dentro de un sistema digital, el código puede conectar datos, herramientas, entornos, permisos y ejecución. Que un modelo domine el coding equivale a que tenga en sus manos la palanca de mando del mundo digital.
\end{knowledgebox}

\subsection{Resumen de la sección}

Coding es un acelerador, no una aplicación aislada. Amplifica a las personas, a las organizaciones y a la propia investigación en IA; por eso podría convertirse en la infraestructura fundamental del segundo acto de la AGI.

